Prompt injection attacks cause LLM agents to execute malicious instructions
embedded in user inputs or tool outputs, overriding system policy.
Existing defenses rely on keyword filters or LLM judges; both operate on
surface text and are either trivially bypassed or themselves
injection-vulnerable.

This thesis investigates whether Abstract Meaning Representation (AMR) can
serve as the basis for static prompt injection detection.
AMR collapses text into canonical predicate-argument graphs, discarding
syntactic variation and surface obfuscation while preserving semantic
intent, making static policy checking tractable.

We propose a scoped formal definition of prompt injection as a payload whose
AMR graph semantically contradicts an explicit policy AMR graph.
The detection mechanism combines static structural rules (polarity
inversion, antonym and synonym predicates) with minimal LLM sub-calls
scoped to single-concept classification queries, a sub-task too narrow for
an attacker to exploit via injection.
Evaluated on a custom attack suite built on the AgentDojo banking benchmark,
the primary configuration achieves high detection rates with zero false
positives, substantially reducing missed attacks over a no-firewall
baseline.
Smatch, an existing AMR graph-similarity metric explored as an alternative,
is rejected because concept-agnostic matching cannot distinguish malicious
instructions from legitimate transaction records.
Remaining misses are tied to AMR parse quality: obfuscation techniques such
as paraphrasing or indirect phrasing can prevent a detection rule from
firing.
